% Sections 5.1 to 5.8, translated from sv/chapters/05/theory.tex.

\section{Powers}
\label{c5:sec:potenser}
A \textbf{power} is a number multiplied by itself repeatedly:
\[
  a^n = \underbrace{a \times a \times \cdots \times a}_{n \text{ factors}}
\]
Here $a$ is called the \textbf{base} and $n$ the \textbf{exponent}.

\begin{exempel}
\[
  2^4 = 2 \times 2 \times 2 \times 2 = 16, \qquad 10^3 = 1\,000
\]
\end{exempel}

\section{Laws of exponents}
\label{c5:sec:potensregler}
The following seven rules hold for all $a \neq 0$ and $b \neq 0$.

\begingroup
\renewcommand{\arraystretch}{1.7}
\begin{narrowtable}{@{}cll@{}}
  \thead{No.} & \thead{Rule} & \thead{Formula} \\
  \midrule
  1 & Product rule & $a^m \cdot a^n = a^{m+n}$ \\
  2 & Quotient rule & $\dfrac{a^m}{a^n} = a^{m-n}$ \\
  3 & Power of a power & $(a^m)^n = a^{m \cdot n}$ \\
  4 & Power of a product & $(a \cdot b)^n = a^n \cdot b^n$ \\
  5 & Power of a quotient & $\left(\dfrac{a}{b}\right)^n = \dfrac{a^n}{b^n}$ \\
  6 & Zero exponent & $a^0 = 1$ \\
  7 & Negative exponent & $a^{-n} = \dfrac{1}{a^n}$ \\
\end{narrowtable}
\endgroup

\begin{exempel}[Rule 1]
\[
  10^3 \cdot 10^4 = 10^{3+4} = 10^7
\]
\end{exempel}

\begin{exempel}[Rule 7]
\[
  10^{-3} = \frac{1}{10^3} = \frac{1}{1000} = 0.001
\]
\end{exempel}

\section{Roots}
\label{c5:sec:rotter}
The \textbf{square root} of $a$ ($a \geq 0$) is the non-negative number whose square is $a$:
\[
  \sqrt{a} = a^{1/2} \quad \text{since} \quad \left(\sqrt{a}\right)^2 = a
\]
The \textbf{$n$th root} of $a$ is the number whose $n$th power is $a$:
\[
  \sqrt[n]{a} = a^{1/n}
\]

\subsection{Fractional exponents}
\label{c5:sec:brakexponenter}
Roots and powers are combined with fractional exponents:
\[
  a^{m/n} = \sqrt[n]{a^m} = \left(\sqrt[n]{a}\right)^m
\]

\begin{exempel}
\[
  8^{2/3} = \left(\sqrt[3]{8}\right)^2 = 2^2 = 4
\]
\end{exempel}

\subsection{Rules for roots}
\label{c5:sec:rotregler}
\[
  \sqrt{a \cdot b} = \sqrt{a} \cdot \sqrt{b}, \qquad \sqrt{\frac{a}{b}} = \frac{\sqrt{a}}{\sqrt{b}}
\]

\begin{aside}
\note[Note] $\sqrt{a + b} \neq \sqrt{a} + \sqrt{b}$. This is a common mistake!
\end{aside}

\section{Scientific notation (standard form)}
\label{c5:sec:standardform}
Electrical engineering involves extremely large and extremely small numbers, for example
$R = 1\,000\,000\,\Omega$ and $C = 0.000\,000\,001\,\text{F}$.

\textbf{Scientific notation} writes the number as $a \times 10^n$, where $1 \leq a < 10$:
\[
  1\,000\,000 = 1.0 \times 10^6, \qquad 0.000\,000\,001 = 1.0 \times 10^{-9}
\]

\subsection{SI prefixes}
\label{c5:sec:siprefix}

\begin{narrowtable}{@{}llc@{}}
  \thead{Prefix} & \thead{Symbol} & \thead{Factor} \\
  \midrule
  Giga & G & $10^9$ \\
  Mega & M & $10^6$ \\
  Kilo & k & $10^3$ \\
  Milli & m & $10^{-3}$ \\
  Micro & $\mu$ & $10^{-6}$ \\
  Nano & n & $10^{-9}$ \\
  Pico & p & $10^{-12}$ \\
\end{narrowtable}

\begin{exempel}
$4.7\,\text{k}\Omega = 4.7 \times 10^3\,\Omega = 4\,700\,\Omega$
\end{exempel}

\begin{exempel}
$33\,\text{nF} = 33 \times 10^{-9}\,\text{F} = 3.3 \times 10^{-8}\,\text{F}$
\end{exempel}

\section{Calculating in scientific notation}
\label{c5:sec:rakning}
\textbf{Multiplication:}
\[
  (3 \times 10^4) \times (2 \times 10^3) = 6 \times 10^7
\]
\textbf{Division:}
\[
  \frac{6 \times 10^6}{2 \times 10^2} = 3 \times 10^4
\]
\textbf{Addition} (the terms must have the same exponent):
\[
  3.5 \times 10^3 + 2.0 \times 10^3 = 5.5 \times 10^3
\]

\section{Significant figures and rounding}
\label{c5:sec:vardesiffror}
\textbf{Significant figures} are the digits that carry meaningful information in a measured value.
\begin{itemize}
  \item $4\,700\,\Omega$ can have 2, 3 or 4 significant figures; there is no telling without
    context.
  \item $4.70 \times 10^3\,\Omega$ clearly has \textbf{3 significant figures}.
\end{itemize}
The result of a calculation should normally be given with as many significant figures as the
least precise input.

\textbf{Rounding rule:} If the digit after the rounding position is $5$ or more, round up;
otherwise round down.

\section{Application: electrical power}
\label{c5:sec:effekt}
The power dissipated in a resistor $R$ at the voltage $U$ is
\[
  P = \frac{U^2}{R}.
\]

\begin{exempel}
With $U = 6\,\text{V}$ and $R = 50\,\Omega$,
\[
  P = \frac{6^2}{50} = \frac{36}{50} = 0.72\,\text{W}.
\]
\end{exempel}

At a given current $I$, the power is
\[
  P = RI^2.
\]

\begin{exempel}
With $I = 20\,\text{mA} = 20 \times 10^{-3}\,\text{A}$ and $R = 1\,\text{k}\Omega = 10^3\,\Omega$,
\[
  P = 10^3 \times (20 \times 10^{-3})^2 = 10^3 \times 400 \times 10^{-6} = 0.4\,\text{W}.
\]
\end{exempel}

\section{Binary and hexadecimal numbers}
\label{c5:sec:binart}
Our everyday number system is a \textbf{positional number system} with base $10$: the value of a
digit depends on where in the number it stands, and each position corresponds to a power of ten.
The rightmost position has the value $10^0 = 1$:
\[
  4\,703 = 4 \cdot 10^3 + 7 \cdot 10^2 + 0 \cdot 10^1 + 3 \cdot 10^0
\]
The same idea works with any base. In digital electronics and programming, the \textbf{binary}
and \textbf{hexadecimal} number systems are used as well as the decimal one. The table shows the
number $45$ in all three:

\begin{narrowtable}{@{}lcll@{\hspace{2em}}l@{}}
  \thead{Number system} & \thead{Base} & \thead{Digits} & \thead{Written} & \thead{In C} \\
  \midrule
  Decimal & $10$ & 0--9 & $45$ & \code{45} \\
  Binary & $2$ & 0 and 1 & $101101_2$ & \code{0b101101} \\
  Hexadecimal & $16$ & 0--9 and A--F & $\mathrm{2D}_{16}$ & \code{0x2D} \\
\end{narrowtable}

The hexadecimal digits A--F stand for $10$--$15$. The small number after a number gives its base;
without one, the number is decimal.

A binary digit is called a \textbf{bit}, and eight bits are called a \textbf{byte}. With $n$ bits,
$2^n$ different values can be stored, from $0$ to $2^n - 1$. A byte can therefore hold $2^8 = 256$
different values, $0$--$255$.

\subsection{From binary or hexadecimal to decimal}
\label{c5:sec:tilldecimalt}
Multiply each digit by the value of its position and add. The positions are numbered from the
right, starting at $0$, so position $n$ has the value $2^n$ in a binary number and $16^n$ in a
hexadecimal number:

\begin{narrowtable}{@{}lcccccccc@{}}
  \thead{Position $n$} & $7$ & $6$ & $5$ & $4$ & $3$ & $2$ & $1$ & $0$ \\
  \midrule
  $2^n$ & $128$ & $64$ & $32$ & $16$ & $8$ & $4$ & $2$ & $1$ \\
\end{narrowtable}

For hexadecimal numbers, $16^0 = 1$, $16^1 = 16$, $16^2 = 256$ and $16^3 = 4\,096$ are usually
enough.

\begin{exempel}
\[
  1011_2 = 1 \cdot 2^3 + 0 \cdot 2^2 + 1 \cdot 2^1 + 1 \cdot 2^0 = 8 + 0 + 2 + 1 = 11
\]
\end{exempel}

\begin{exempel}
F stands for $15$:
\[
  \mathrm{2F}_{16} = 2 \cdot 16^1 + 15 \cdot 16^0 = 32 + 15 = 47
\]
\end{exempel}

\subsection{From decimal to binary or hexadecimal}
\label{c5:sec:frandecimalt}
Divide the number by the base repeatedly and note the remainder each time, until the quotient is
$0$. The remainders, read \textbf{from bottom to top}, are the number in the new base.

\begin{exempel}[From decimal to binary]
Convert $13$ to binary.

\begin{narrowtable}{@{}ccc@{}}
  \thead{Division} & \thead{Quotient} & \thead{Remainder} \\
  \midrule
  $13 / 2$ & $6$ & $1$ \\
  $6 / 2$ & $3$ & $0$ \\
  $3 / 2$ & $1$ & $1$ \\
  $1 / 2$ & $0$ & $1$ \\
\end{narrowtable}

The remainders from bottom to top give $13 = 1101_2$. Check: $8 + 4 + 0 + 1 = 13$~\checkmark
\end{exempel}

\begin{exempel}[From decimal to hexadecimal]
Convert $470$ to hexadecimal. A remainder between $10$ and $15$ is written as a letter.

\begin{narrowtable}{@{}ccc@{}}
  \thead{Division} & \thead{Quotient} & \thead{Remainder} \\
  \midrule
  $470 / 16$ & $29$ & $6$ \\
  $29 / 16$ & $1$ & $13 = \mathrm{D}$ \\
  $1 / 16$ & $0$ & $1$ \\
\end{narrowtable}

The remainders from bottom to top give $470 = \mathrm{1D6}_{16}$. Check:
$1 \cdot 256 + 13 \cdot 16 + 6 = 470$~\checkmark
\end{exempel}

\begin{aside}
\note[Tip] A calculator gives the quotient as a decimal number. $470 / 16 = 29.375$ means the
quotient $29$ and the remainder $0.375 \cdot 16 = 6$.
\end{aside}

\subsection{Between binary and hexadecimal}
\label{c5:sec:binarthex}
Since $16 = 2^4$, each hexadecimal digit corresponds to exactly four bits:

\begin{narrowtable}{@{}cc@{\hspace{2em}}cc@{\hspace{2em}}cc@{\hspace{2em}}cc@{}}
  \thead{Hex} & \thead{Binary} & \thead{Hex} & \thead{Binary} &
  \thead{Hex} & \thead{Binary} & \thead{Hex} & \thead{Binary} \\
  \midrule
  $0$ & $0000$ & $4$ & $0100$ & $8$ & $1000$ & $\mathrm{C}$ & $1100$ \\
  $1$ & $0001$ & $5$ & $0101$ & $9$ & $1001$ & $\mathrm{D}$ & $1101$ \\
  $2$ & $0010$ & $6$ & $0110$ & $\mathrm{A}$ & $1010$ & $\mathrm{E}$ & $1110$ \\
  $3$ & $0011$ & $7$ & $0111$ & $\mathrm{B}$ & $1011$ & $\mathrm{F}$ & $1111$ \\
\end{narrowtable}

\textbf{From binary to hexadecimal:} split the bits into groups of four, \textbf{from the right},
pad with zeros on the left if needed, and replace each group with its hexadecimal digit:
\[
  1011\,0110_2 = \underbrace{1011}_{\mathrm{B}}\;\underbrace{0110}_{6} = \mathrm{B6}_{16}
\]
\textbf{From hexadecimal to binary:} replace each digit with its four bits:
\[
  \mathrm{4D}_{16} = \underbrace{0100}_{4}\;\underbrace{1101}_{\mathrm{D}} = 0100\,1101_2
\]
A byte therefore corresponds to exactly two hexadecimal digits, from \code{0x00} to \code{0xFF}.
This is why, for example, the contents of a microcontroller's registers are often written in
hexadecimal: \code{0xB6} is easier to read than \code{0b10110110}.

\begin{aside}
\note[Note] In computing, kilo is sometimes used for $2^{10} = 1\,024$ instead of $1\,000$. The
unambiguous name is \textbf{kibi} (Ki): $1\,\text{KiB} = 1\,024$~bytes.
\end{aside}
